\chapter[Программная реализация детектора временной ошибки(синхронизация приемника и передатчика) на SDR]{Программная реализация детектора временной ошибки(синхронизация приемника и передатчика) на SDR. Написание функций петли (контура) синхронизации}

\begin{leftbar}
    \href{https://github.com/someotb/sdr_reports/tree/main/practices/practice7}{Ссылка на GitHub}
\end{leftbar}

\section*{Цель практики:}

Написать функцию которая будет автоматически синхронизировать примник и передатчик.

\section*{Краткие теоретические сведения}

Оптимальный приемник выполняет прием дискретных сигналов с минимальной вероятностью ошибки. Строится оптимальный приемник по схеме на
основе корелляторов или согласованных фильтров.

\vspace{10pt}

Оптимальный приемник выбирает наиболее близкое принятому сигналу r(t)
из набора возможных колебаний si(t). Это может выполняться вычислением
расстояния

\begin{figure}[H]
    \centering
    \includegraphics[width=1.0\textwidth]{mf.png}
    \caption{Matched filter}    
\end{figure}

Если использовать формирующий фильтр с прямоугольной формой ИХ, то
форма ИХ согласованного фильтра также является прямоугольной. Сигнал на
выходе СФ вычисляется как свертка отсчетов входного сигнала и ИХ согласованного фильтра.

\begin{figure}[H]
    \centering
    \includegraphics[width=1.0\textwidth]{mf_cor.png}
    \caption{Matched filtet/correlator output}    
\end{figure}

Демодулятор должен сделать отсчет сигнала на выходе СФ точно в момент
максимума сигнала на выходе СФ. Если время взятия отсчета не совпадает с
максимумом, то невозможно принять надежное решение о передаваемом символе.

\begin{figure}[H]
    \centering
    \includegraphics[width=1.0\textwidth]{perfect_bpsk.png}
    \caption{Идеальный отсчет}    
\end{figure}

Время взятия отсчета (symbol timing) совпадает с максимумом сигнала на выходе СФ. При смещении времени взятия отсчета берутся значения с выхода СФ
не в оптимальные моменты времени, что приводит к нарушению возможности
точного определения амплитуды символа.

\begin{figure}[H]
    \centering
    \includegraphics[width=1.0\textwidth]{verybad_meh.png}
    \caption{Плохой отсчет}    
\end{figure}

\textbf{Детектор временной ошибки - TED}\newline
Для поиска нужного отсчета сначала вычисляется ошибка положения символа при помощи детектора временной ошибки. Детектор использует набор отсчетов и вычисляет ошибку ТЕД e(n), величина которой зависит от смещения
максимума выхода СФ. В таблице приведены выражения для вычисления величины ошибки по отсчетам сигнала с выходы СФ.

\begin{figure}[H]
    \centering
    \includegraphics[width=1.0\textwidth]{ted.png}
    \caption{TED}    
\end{figure}

Зависимость сигнала на выходе TED от величины смещения символа (временной ошибки) называется S-curve, для детектора Гарднера она имеет форму

\begin{figure}[H]
    \centering
    \includegraphics[width=0.7\textwidth]{ted_graph.png}
    \caption{Форма TED}    
\end{figure}

\textbf{Фильтр контура}\newline
Фильтр выполнен в виде пропорционально-интегрирующего звена по схеме

\begin{figure}[H]
    \centering
    \includegraphics[width=0.7\textwidth]{count_filter.png}
    \caption{Форма TED}    
\end{figure}

Параметры фильтра

\begin{figure}[H]
    \centering
    \includegraphics[width=0.7\textwidth]{params.png}
    \caption{Параметры фильтра}    
\end{figure}
 
\indent Программно фильтр реализуется в общем цикле по отсчетам сигнала на
выходе СФ.\\
\indent Начальные значения переменных в ветвях K1 и K2 - p1 и p2 равны 0. Сами
величины K1 и K2 вычисляются перед циклом. Нормированная полоса фильтра
контура BnTs = 0.01. Также перед циклом задается начальное смещение offset
=0 (номер отсчета символа, один из 10).\\
\indent В цикле по отсчетам сигнала на выходе СФ вычисляется текущее значение
ошибки e для текущего смещения n = offset;\\
\indent e = (I(n+Nsp+Ns) - I(n+Ns)) * I(n+Nsp/2+Ns) + (Q(n+Nsp+Ns) - Q(n+Ns))
* Q(n+Nsp/2+Ns).\\
\indent Если e - текущая величина ошибки на выходе детектора ошибки (на данном
индексе отсчета сигнала на выходе СФ), то в ветви K1 переменная p1 = e*K1.
В ветви K2 переменная p2 = p2 + p1 + e*K2. Эта переменая p2 имеет смысл
задержки, дробного смещения (offset) на длительности символа. Тут же , если
p2>1, ее сбрасывают\\
\indent while(p2 > 1) p2 = p2 - 1; end.\\
\indent После этого определяется уточнененное смещение в отсчетах (offset), для
этого нужно умножить p2 на Nsps и округлить произведение до ближайшего
целого, offset = round(p2*Nsps), это будет смещение на данной итерации (оптимальный отсчет текущего символа, индекс в массиве отсчетов сигнала на выходе СФ).
На следующем символе снова вычисляется
величина ошибки и после фильтрации определяется смещение, по которому находится следующий индекс отсчета offset+Nsps . Таким образом контур сводит
ошибку к мимнимуму и определяются оптимальные индексы в массиве отсчетов
символов сигнала на выходе СФ.
\section*{Выполнение}

Напишем функцию Гарднера, которая будет автоматически искать offset

\begin{lstlisting}[language=python, caption=Функция Гарднера]
import numpy as np
def gardner(complex_symbols_saw):
    p1 = 0
    p2 = 0
    offset = 0
    BnTs = 0.01
    Nsp = 10 # Samples per symbol
    zeta = np.sqrt(2)/2
    teta = (BnTs/Nsp)/(zeta + 1/(4*zeta))
    css = complex_symbols_saw
    Kp = 0.002
    K1 = (-4 * zeta * teta)/((1 + 2 * zeta * teta + teta ** 2) * Kp)
    K2 = (-4 * teta ** 2)/((1 + 2 * zeta * teta + teta ** 2) * Kp)
    list_of_offset = []
    
    for n in range(0, len(css)-22):
        err = (np.real(css[n + Nsp + Nsp]) - np.real(css[n + Nsp])) * np.real(css[n + Nsp + Nsp//2]) + (np.imag(css[n + Nsp + Nsp]) - np.imag(css[n + Nsp])) * np.imag(css[n + Nsp + Nsp//2])
        err = np.real(err)
        p1 = err * K1
        p2 = p2 + p1 + err * K2
        while p2 > 1:
            p2 = p2 - 1
        while p2 < -1:
            p2 = p2 + 1
        offset = np.round(p2 * Nsp)
        list_of_offset.append(offset)
        
    return int(offset), list_of_offset
\end{lstlisting}

\vspace{20pt}

Используем эту функций в нашем коде

\begin{lstlisting}[language=python, caption=Функция Гарднера]
from matplotlib import pyplot as plt
import numpy as np
from auto_offset import gardner

rx = np.fromfile(f"/home/someotb/Files/Code/sdr/pluto/dev/build/rx.pcm", dtype=np.int16)

samples = []

for x in range(0, len(rx), 2):
    samples.append((rx[x] + 1j * rx[x+1])/np.max(rx))

plt.figure()
plt.scatter(np.real(samples), np.imag(samples))
plt.axhline()
plt.axvline()

impulse = np.ones(10)

I = np.real(samples)
Q = np.imag(samples)

i_convolve = np.convolve(I, impulse)
q_convolve = np.convolve(Q, impulse)

i_con_dec = []
q_con_dec = []
list_offset = []

offset, list_offset = gardner(i_convolve + 1j * q_convolve)
print(offset)

plt.figure()
plt.title("offset")
plt.plot(np.arange(len(list_offset)), list_offset)
plt.ylim([-20, 20])
plt.xlim([4060, 4560])

for i in range(offset, len(i_convolve), 10):
    i_con_dec.append(i_convolve[i])
    q_con_dec.append(q_convolve[i])

# Убираем шумы
# i_con_dec = [x for x in i_con_dec if not (abs(x) <= 1)]
# q_con_dec = [x for x in q_con_dec if not (abs(x) <= 1)]

print("lenght i_con_dec: ", len(i_con_dec))
print("lenght q_con_dec: ", len(q_con_dec))

plt.figure()
plt.subplot(2,1,1)
plt.scatter(i_con_dec, q_con_dec)
plt.axhline()
plt.axvline()
plt.subplot(2,1,2)
plt.xlim([5000, 5500])
plt.ylim([-30, 30])
plt.scatter(np.arange(len(i_convolve)), i_convolve)
plt.scatter(np.arange(len(q_convolve)), q_convolve)

plt.show()
\end{lstlisting}

\section*{Вывод}
Написал, программную реализация детектора ошибки, исполнил в коде схему Гарднера.